\editfromrepo{base}

\replace{What is the difference between these two sections? Why are both
necessary?}{%
\label{n:contract-deed}%
As indicated by this statute, real estate sales typically require not one but
two formal documents. The \term{contract of sale} is a promise between the
parties to make the transaction: The seller agrees to convey the land, and the
buyer agrees to pay the negotiated price. The \term{deed}, by contrast, is the
document that actually causes property rights in the land to transfer from buyer
to seller. (The deed does the deed, indeed.)
\endgraf
Nothing stops these two documents from being executed at the same time, or even
on the same piece of paper. Yet typically when buying and selling land, there is
an \term{executory period} of typically a full month between signing the
contract and \term{closing}, when the deed is executed. What's going on during
that month? In particular, why might either of the parties to the contract want
to back out of the sale, and what contingencies might they insert into the
contract? For some ideas, look up any guide on buying and selling a home. }

\endedit
